% Section 7.6, from lectures/L07/README.md.

\section{Summary}
\label{c7:sec:review}

\needspace{6\baselineskip}
\subsection*{What you should be able to do now}
\begin{itemize}
  \item Explain what separates a timer from an ordinary counter, and why the difference is a
    comparison.
  \item Build a timer as a gate network: an equality comparator of one XNOR per bit, and the clear
    that makes it repeat instead of wrapping.
  \item Implement a timer in VHDL as a reusable module, convert a period in seconds into a
    \code{TICK_COUNT}, and say why \code{timeout} is a one-clock-cycle pulse rather than a level.
  \item Compose a design out of finished modules instead of rewriting them, and say why a button
    controlling a timer still has to pass through a synchroniser first.
\end{itemize}

\needspace{6\baselineskip}
\subsection*{Questions to test yourself with}
\begin{enumerate}
  \item In what sense is a timer ``just a counter with a comparison''?
  \item Why must the internal counter be cleared to \code{0} when it reaches its target, instead of
    being left to count on?
  \item Why is \code{timeout} a single-cycle pulse rather than a level that stays high?
  \item In the CircuitVerse version, changing the target value means rewiring the comparator. What
    replaces that rewiring in the VHDL version, and when is its value fixed?
  \item Why is it still necessary to run a button through a synchroniser before its edge is used to
    enable a timer, even though the button has nothing to do with the counting?
  \item \code{walking_led} writes almost no logic of its own. Name the three modules it composes and
    which chapter each came from.
\end{enumerate}

\needspace{6\baselineskip}
\subsection*{Looking ahead}
\Chapref{c8:ch} goes on with:
\begin{itemize}
  \item State machines, the book's last building block, first designed by hand: state diagram,
    state table, Karnaugh-derived logic and a gate network in CircuitVerse, exactly as this
    chapter's timer was drawn before it was written.
  \item Then in VHDL, where an enumerated type and a \code{case} replace the state table and the
    Karnaugh maps entirely.
  \item A timer is itself a state machine with two states and a single purpose. \Chapref{c8:ch}
    generalises that to circuits with many states and many transitions between them.
\end{itemize}
